\section{Discussion}\label{sec:discussion}

\paragraph{Reconstruction is a proxy for behaviour, not for information.} Within a layer, an SAE with lower
reconstruction error damages the model's predictions less, but it does not keep more truth information. Reconstruction
metrics average over all directions weighted by their variance; the information that distinguishes true from false
statements, or honest from dishonest instructions, lies in small differences between nearly identical activations.

\paragraph{Why the causal effect leaks.} M4 makes this concrete. An SAE can reconstruct the honest and the dishonest
activation well individually and still reconstruct their difference poorly, because the two reconstruction errors do
not cancel. At mid and late sites about a third of the effect is outside the reconstructed difference, and for TopK
and JumpReLU part of it exists only when reconstruction and error act together through GPT-2's nonlinear downstream
computation.

\paragraph{What fine-tuning changes, and what it does not.} The fine-tuning experiments separate three things
that are often reported as one. (i) \emph{Task reconstruction} improves reliably: reconstruction error on truth
statements falls by 61--85\,\% in every cell. (ii) \emph{Frozen-probe preservation} improves in direction
everywhere and meets the pre-registered criterion in 4 of 9 cells, but a re-trained reader recovers no more truth
information from the fine-tuned reconstruction than from the original one (no cell changes consistently). Because M2
showed that frozen-probe gaps are dominated by the probe's sensitivity to off-distribution perturbations, the most
economical reading is that reconstruction fine-tuning moves the reconstruction back towards the region where the
original probe's decision boundary was fitted, rather than restoring information the SAE had discarded. (iii)
\emph{Behavioural preservation} improves with an output-KL term in all cells, and creditably in 5 of 9; it does not
transfer to truth preservation, which is unchanged or slightly worse. Improving one level of faithfulness therefore
does not improve the others, and a claim of ``more faithful SAEs'' must name the level.

\paragraph{Sparsity is part of the price.} For ReLU and JumpReLU at layers 8 and 9 the output-KL gain comes with a
5--8\,\% rise in L0 that the continued sparsity controller does not correct. Since lower reconstruction error and
lower KL are both easier at higher L0, behavioural gains reported without a sparsity guard would overstate the
benefit; TopK, whose L0 is fixed, is the only architecture where the gain is creditable at every layer.

\paragraph{The negative result is the main finding.} The SAEs studied here reconstruct well by conventional measures,
yet they lose task information that PCA at equal error keeps, and at mid and late layers their reconstruction carries
only part of a causal effect that GPT-2 demonstrably carries there. Short fine-tuning improves what it optimises
(task reconstruction, output agreement) and nothing we could verify beyond that. Whether fine-tuning also improves
\emph{causal} preservation is the open question addressed by E3. All of these are \evlabel{validation} findings.

\pending{Revise after E1--E3 and M5; in particular, whether the conclusions hold for an externally trained SAE
(E2) and at task-level matched sparsity (E1).}
